\section{Introduction}
\label{intro}

Supervised learning remains notoriously weak at generalization under distribution shifts.
Unless training and test data are drawn from the same distribution, even seemingly minor differences turn out to defeat state-of-the-art models \citep{recht2018cifar}.
Adversarial robustness and domain adaptation are but a few existing paradigms that try to \textit{anticipate} differences between the training and test distribution with either topological structure or data from the test distribution available during training.
We explore a new take on generalization that \textit{does not} anticipate the distribution shifts, but instead learns from them at test time.

We start from a simple observation. 
The unlabeled test sample $x$ presented at test time gives us a hint about the distribution from which it was drawn. 
We propose to take advantage of this hint on the test distribution by allowing the model parameters $\vtheta$
to depend on the test sample $x$, but not its unknown label $y$.
The concept of a variable decision boundary $\vtheta(x)$ is powerful in theory since it breaks away from the limitation of fixed model capacity (see additional discussion in \Cref{vdb}), but the design of a feedback mechanism from $x$ to $\vtheta(x)$ raises new challenges in practice that we only begin to address here.

Our proposed test-time training method creates a self-supervised learning problem based on this single test sample $x$, updating $\vtheta$ at test time before making a prediction. 
Self-supervised learning uses an auxiliary task that automatically creates labels from unlabeled inputs.
In our experiments, we use the task of rotating each input image by a multiple of 90 degrees and predicting its angle~\citep{gidaris2018unsupervised}.

This approach can also be easily modified to work outside the standard supervised learning setting. 
If several test samples arrive in a batch, we can use the entire batch for test-time training.
If samples arrive in an online stream, we obtain further improvements by keeping the state of the parameters.
After all, prediction is rarely a single event. The online version can be the natural mode of deployment under the additional assumption that test samples are produced by the same or smoothly changing distribution shifts.

We experimentally validate our method in the context of object recognition on several standard benchmarks. These include images with diverse types of corruption at various levels \citep{hendrycks2019benchmarking},  video frames of moving objects \citep{shankar2019systematic}, and a new test set of unknown shifts collected by \cite{recht2018cifar}. 
Our algorithm makes substantial improvements under distribution shifts, while maintaining the same performance on the original distribution.

In our experiments, we compare with a strong baseline (labeled joint training) that uses both supervised and self-supervised learning at training-time, but keeps the model fixed at test time.
Recent work shows that \emph{training-time} self-supervision improves robustness \citep{hendrycks2019using}; our joint training baseline corresponds to an improved implementation of this work. 
A comprehensive review of related work follows in \Cref{related}.

We complement the empirical results with theoretical investigations in \Cref{theory}, and establish an intuitive sufficient condition on a convex model of when Test-Time Training helps; this condition, roughly speaking, is to have correlated gradients between the loss functions of the two tasks.
\newcommand\blfootnote[1]{%
  \begingroup
  \renewcommand\thefootnote{}\footnote{#1}%
  \addtocounter{footnote}{-1}%
  \endgroup
}
\blfootnote{Project website: {\scriptsize \url{https://test-time-training.github.io/}}.}
\newpage